\newcommand{\detourpanel}[4]{%
\begin{picture}(156,86)
  \put(14,26){\circle{19}}\put(14,26){\makebox(0,0){#1}}
  \put(74,26){\circle{19}}\put(74,26){\makebox(0,0){#2}}
  \put(134,26){\circle{19}}\put(134,26){\makebox(0,0){#4}}
  \put(74,66){\circle{19}}\put(74,66){\makebox(0,0){#3}}
  \put(24,26){\vector(1,0){40}}\put(84,26){\vector(1,0){40}}
  \put(70,35){\vector(0,1){22}}\put(78,57){\vector(0,-1){22}}
  \put(104,36){\makebox(0,0){\scriptsize locked}}
  \put(74,82){\makebox(0,0){\scriptsize key}}
  \put(14,7){\makebox(0,0){\scriptsize start}}
  \put(134,7){\makebox(0,0){\scriptsize goal}}
\end{picture}}
\begin{figure}[H]
\centering
\setlength{\unitlength}{1pt}
\begin{tabular}{cc}
\detourpanel{A}{D}{C}{B} & \detourpanel{F}{C}{E}{A} \\
Original labels & Equivalent labels \\
$A\to D\to C\to D\to B$ & $F\to C\to E\to C\to A$
\end{tabular}
\caption{Schematic key detour and an equivalent relabeling. The two
paths describe the same move sequence after mapping the vertex names.
Extra graph vertices and distractor edges are omitted from the schematic.}
\label{fig:detour}
\end{figure}
